\documentclass[11pt]{article}

% ---- Portable preamble (compiles with XeLaTeX or pdfLaTeX) ----
% Notation is kept deliberately consistent with lemma1-bias-proof.tex so this
% note can be imported into the manuscript alongside Lemma 1 without renaming.
\usepackage{amsmath,amsthm,amssymb}
\usepackage[margin=1in]{geometry}
\usepackage[colorlinks=true,linkcolor=blue,urlcolor=blue,citecolor=blue]{hyperref}
\usepackage{mathtools}
\usepackage[numbers,round]{natbib}

% ---- Theorem environments ----
\theoremstyle{plain}
\newtheorem{theorem}{Theorem}
\newtheorem{lemma}{Lemma}
\newtheorem{proposition}{Proposition}
\newtheorem{corollary}{Corollary}
\theoremstyle{definition}
\newtheorem{assumption}{Assumption}
\newtheorem{definition}{Definition}
\theoremstyle{remark}
\newtheorem*{remark}{Remark}

% ---- Notation shortcuts (match lemma1-bias-proof.tex) ----
\newcommand{\E}{\mathbb{E}}
\newcommand{\PP}{\mathbb{P}}
\newcommand{\Real}{\mathbb{R}}
\newcommand{\indep}{\perp\!\!\!\perp}
\newcommand{\given}{\,\vert\,}
\newcommand{\ef}{e_f}
\newcommand{\ez}{e_0}
\newcommand{\mz}{m_0}
\newcommand{\fz}{f_0}
\newcommand{\Ltwo}[1]{\lVert #1 \rVert_{2}}
\DeclarePairedDelimiter{\abs}{\lvert}{\rvert}
\DeclarePairedDelimiter{\norm}{\lVert}{\rVert}
\newcommand{\Pn}{\mathbb{P}_n}
\newcommand{\eT}{e_{\hat\tau}}
\newcommand{\mT}{m_{\hat\tau}}
\newcommand{\etau}{e_{\tau_0}}
\newcommand{\mtau}{m_{\tau_0}}
\newcommand{\thstar}{\theta^{\star}_n}
\newcommand{\htheta}{\hat\theta}
\newcommand{\psiscore}{\psi}
\newcommand{\Ghat}{\hat g}
\newcommand{\gz}{g_0}
\newcommand{\odiff}{o}

\title{Working Note: Single Interpretable Tree per Nuisance \\
  Valid Full-Sample Inference for the Working-Model ATT via LOO Stability}
\author{}
\date{2026-08-04}

\begin{document}
\maketitle

\begin{abstract}
We prove that a \emph{single} interpretable decision tree per nuisance, with a
\emph{fixed} bounded-complexity partition and leaf values refit on \emph{all} $n$
observations, yields valid $\sqrt{n}$ inference for the ATT of the displayed tree
model---using leave-one-out (LOO) algorithmic stability
\citep{ChenSyrgkanisAustern2022} rather than sample-splitting or cross-fitting.
The target is the \emph{working-model ATT} $\thstar$ (Lemma~1); the true ATT
$\theta_0$ is generically \emph{not} targetable at $\sqrt{n}$ by any
fixed-complexity interpretable tree (Corollary~\ref{cor:impossible}), and the
gap $\thstar-\theta_0$ is reported as an honest diagnostic $\hat\delta$, never
used to correct the estimate. A structural piecewise-constant assumption
recovers $\theta_0$ as a special case (Corollary~\ref{cor:structural}).
\end{abstract}

\section{Setup and notation}\label{sec:setup}

Let $O=(X,A,Y)\sim P$ be i.i.d.\ with covariates $X\in[0,1]^d$, binary
treatment $A\in\{0,1\}$, and outcome $Y$. Write $\pi=\PP(A=1)>0$,
$\ez(X)=\PP(A=1\given X)$ (propensity), and $\mz(X)=\E[Y\given X,A=0]$ (control
outcome regression). The estimand is the average treatment effect on the treated
$\theta_0=\E[Y(1)-Y(0)\given A=1]$, identified under unconfoundedness
$\{Y(0),Y(1)\}\indep A\given X$ and positivity (Assumption~\ref{ass:positivity}).

The Neyman-orthogonal ATT score, in the empirical treated-mean identification
(no $m_1$ regression), is (matching eq.~(score) of the Lemma~1 note)
\begin{equation}\label{eq:score}
  \psiscore(O;\theta,e,f)
  = \frac{A}{\pi}\bigl(Y-f(X)-\theta\bigr)
    - \frac{1}{\pi}\,\frac{e(X)\,(1-A)}{1-e(X)}\,\bigl(Y-f(X)\bigr),
\end{equation}
where $(e,f)$ are candidate nuisances plugged in for $(\ez,\mz)$. For any
$(e,f)\in L^2(P)$ the \emph{working-model ATT} $\theta(e,f)$ solves
$\E_P[\psiscore(O;\theta,e,f)]=0$; at the truth $\theta(\ez,\mz)=\theta_0$.
Lemma~1 (imported) gives the bilinear, doubly robust bias
\begin{equation}\label{eq:lemma1}
  \theta(e,f)-\theta_0
  = \frac{1}{\pi}\,\E_P\!\left[\frac{\ez-e}{1-e}\,(\mz-f)\right],
  \qquad
  \abs{\theta(e,f)-\theta_0}\le \frac{1}{\pi c}\,\Ltwo{e-\ez}\,\Ltwo{f-\mz},
\end{equation}
under Assumption~\ref{ass:positivity} (the constant is $1/(\pi c)$; see the
Lemma~1 note's verification remark).

\paragraph{Trees and the working-model estimand.}
A partition $\tau$ of $[0,1]^d$ into $L=\abs{\tau}$ axis-aligned cells (leaves)
induces piecewise-constant nuisances by cell-wise population means. On the
project's adaptive grid $G_n$ (quantile bins $b_n\asymp\log n$ per feature; see
the discretization note), let $\tau_0=\tau_0(G_n)$ be the \emph{grid-optimal}
partition of bounded complexity $L=\odiff(\sqrt n)$, with piecewise-constant
nuisances $\etau,\mtau$. Define the \textbf{interpretable working-model ATT}
\begin{equation}\label{eq:thetastar}
  \thstar \;:=\; \theta(\etau,\mtau),
\end{equation}
the ATT functional \eqref{eq:score} evaluated at the grid-optimal tree model.
This---not $\theta_0$---is the target of the inference below.

\paragraph{The estimator.}
Given a data-selected partition $\hat\tau$ (from cross-fit trees or a
Rashomon/modal menu), fix it and refit leaf values as within-leaf sample means
on all $n$: $\eT(x)=\Pn[A\mid X\in\ell(x)]$,
$\mT(x)=\Pn[Y\mid X\in\ell(x),A=0]$, where $\ell(x)$ is the leaf containing $x$.
The full-sample plug-in $\htheta$ solves $\Pn[\psiscore(O;\theta,\eT,\mT)]=0$.

\paragraph{CSA linear-moment form.}
Write the moment in the linear form of \citet{ChenSyrgkanisAustern2022},
$m(O;\theta,g)=a(O;g)\,\theta+\nu(O;g)$ with nuisance $g=(e,f)$. From
\eqref{eq:score},
\begin{equation}\label{eq:anu}
  a(O;g)=-\frac{A}{\pi},\qquad
  \nu(O;g)=\frac{A}{\pi}\bigl(Y-f(X)\bigr)
    -\frac{1}{\pi}\frac{e(X)(1-A)}{1-e(X)}\bigl(Y-f(X)\bigr),
\end{equation}
so $a$ does not depend on $g$ (hence $A(g)\equiv-1/\pi\cdot\pi=-1$ after taking
expectations, giving $A(g_0)^{-1}=-1$) and all nuisance dependence sits in
$\nu$. This simplifies the CSA conditions: only the $\nu$-stability
conditions (S3)--(S4) are nonvacuous.

\begin{assumption}[Positivity]\label{ass:positivity}
There is $c\in(0,\tfrac12]$ with $\ez(X)\in[c,1-c]$ a.s., and every candidate
tree propensity satisfies $\eT(X),\etau(X)\in[c,1-c]$ (enforced by leaf-level
clipping / a minimum leaf treated-fraction).
\end{assumption}

\begin{assumption}[Bounded outcome moments]\label{ass:moments}
$\E[Y^2\mid X]\le C_Y<\infty$ a.s.
\end{assumption}

\begin{assumption}[Balanced bounded complexity]\label{ass:leaves}
$\hat\tau$ has $L=\abs{\hat\tau}$ leaves with $L=\odiff(\sqrt n)$ and each leaf
count $n_\ell\ge \underline{c}\,n/L$ for a constant $\underline{c}>0$
(min-leaf-size regularization). Hence $\min_\ell n_\ell\gtrsim n/L\gg\sqrt n$.
\end{assumption}

\section{Leaf-refit LOO stability}\label{sec:leaf}

We first show the leaf-refit map is LOO-stable at the CSA rate for a
\emph{fixed} partition. Let $\Ghat=(\eT,\mT)$ and let $\Ghat^{(-l)}$ be the same
map with observation $O_l$ replaced by an independent copy $\tilde O_l$ (the
partition $\hat\tau$ held fixed). Recall the CSA stability quantities for the
$\nu$-component: with $Z=O$ a fresh point,
\begin{align}
  \text{(S3)}\quad & \max_{l\le n}\bigl\lVert \nu(O_l;\Ghat)-\nu(O_l;\Ghat^{(-l)})\bigr\rVert_1=\odiff(n^{-1/2}),\label{eq:S3}\\
  \text{(S4)}\quad & \max_{l\le n}\bigl\lVert \nu(O;\Ghat)-\nu(O;\Ghat^{(-l)})\bigr\rVert_2=\odiff(n^{-1/2}),\label{eq:S4}
\end{align}
and the mean-squared-continuity (MSC) condition: for some $q<\infty$, $L_{\mathrm{MSC}}>0$,
$\E[(\nu(O;g)-\nu(O;g'))^2]\le L_{\mathrm{MSC}}\Ltwo{g-g'}^{q}$ for all $g,g'$.

\begin{lemma}[Leaf-refit LOO stability]\label{lem:leafstab}
Under Assumptions~\ref{ass:positivity}--\ref{ass:leaves} and a \emph{fixed}
partition $\hat\tau$, the leaf-mean map $\Ghat=(\eT,\mT)$ satisfies the CSA
$\nu$-stability conditions \eqref{eq:S3}--\eqref{eq:S4} and the MSC condition with
$q=2$. Consequently the CSA stochastic-equicontinuity condition holds for the
moment \eqref{eq:anu}.
\end{lemma}

\begin{proof}
\emph{Step 1: MSC with $q=2$.} On the positivity set $e\in[c,1-c]$, the map
$e\mapsto e/(1-e)$ is Lipschitz with constant $\le 1/c^2$ (its derivative is
$1/(1-e)^2\le 1/c^2$). Hence, writing $g=(e,f)$, $g'=(e',f')$, from
\eqref{eq:anu} and $\abs{Y-f}\le\abs{Y}+\abs{f}$ with
Assumption~\ref{ass:moments},
\[
  \abs{\nu(O;g)-\nu(O;g')}
  \le \frac{1}{\pi}\abs{f-f'}
    +\frac{1}{\pi}\Bigl(\tfrac{1}{c^2}\abs{e-e'}\abs{Y-f}
      +\tfrac{1-c}{c}\abs{f-f'}\Bigr),
\]
so squaring, taking expectations, and using boundedness gives
$\E[(\nu(O;g)-\nu(O;g'))^2]\le L_{\mathrm{MSC}}\bigl(\Ltwo{e-e'}^2+\Ltwo{f-f'}^2\bigr)
=L_{\mathrm{MSC}}\Ltwo{g-g'}^2$, i.e.\ $q=2$.

\emph{Step 2: LOO perturbation of a leaf mean is $O(1/n_\ell)$.} Fix a leaf
$\ell$ with count $n_\ell$ and control-arm count $n_\ell^0$. Replacing one
observation $O_l$ changes the within-leaf mean by at most
\[
  \abs{\mT^{(-l)}(x)-\mT(x)}
  \le \frac{\abs{Y_l}+\abs{\tilde Y_l}}{n_\ell^0}
  \quad\text{if }X_l\in\ell,\ A_l=0,\qquad 0\text{ otherwise,}
\]
and analogously $\abs{\eT^{(-l)}(x)-\eT(x)}\le 1/n_\ell$. By
Assumption~\ref{ass:leaves}, $n_\ell,n_\ell^0\gtrsim n/L$, so under
Assumption~\ref{ass:moments},
$\Ltwo{\Ghat-\Ghat^{(-l)}}=O(L/n)$ uniformly in $l$.

\emph{Step 3: (S3)--(S4).} By the MSC bound (Step 1, $q=2$) and Step 2, for a
fresh $O$,
\[
  \bigl\lVert \nu(O;\Ghat)-\nu(O;\Ghat^{(-l)})\bigr\rVert_2
  \le L_{\mathrm{MSC}}^{1/2}\,\Ltwo{\Ghat-\Ghat^{(-l)}}
  = O(L/n).
\]
The $\lVert\cdot\rVert_1$ version at the training point $O_l$ is bounded the
same way (a single point contributes $O(1/n_\ell)=O(L/n)$ to the mean, and
$\lVert\cdot\rVert_1\le\lVert\cdot\rVert_2$). Since $L=\odiff(\sqrt n)$,
$O(L/n)=\odiff(n^{-1/2})$, giving \eqref{eq:S3}--\eqref{eq:S4}.

\emph{Step 4.} By \citet[Lemma~2]{ChenSyrgkanisAustern2022}, (S3)--(S4)+MSC imply
the stochastic-equicontinuity condition for \eqref{eq:anu} (the $a$-component is
$g$-free, so (S1)--(S2) are vacuous).
\end{proof}

\section{Main result: full-sample CLT for the working-model ATT}\label{sec:main}

Selection instability is quarantined by a \emph{partition-fixing event} $E_n$,
under either of two mechanisms.

\begin{assumption}[Partition-fixing event]\label{ass:En}
Let $E_n=\{\hat\tau=\tau_0\}$ be the event that the selector recovers the
grid-optimal partition. We require the failure probability to vanish
\emph{faster than the root-$n$ rate},
\begin{equation}\label{eq:En-rate}
  \PP(E_n^c)\;=\;\odiff\!\bigl(n^{-1/2}\bigr).
\end{equation}
This is strictly stronger than $\PP(E_n)\to1$ and is the operative requirement:
$\odiff(1)$ does \emph{not} suffice (see the proof of Theorem~\ref{thm:clt}, final
step, and Remark~\ref{rem:rate-necessity}). Condition \eqref{eq:En-rate} holds
under either mechanism:
\begin{itemize}
  \item[(B)] \emph{Margin / exact recovery.} On the grid $G_n$ the grid-optimal
  $\tau_0$ has a risk margin $\Delta_n>0$ over its competitors, and a modal
  selector over $M_n$ independent splits obeys a union-bound tail
  $\PP(E_n^c)\le S_n\exp(-M_n\Delta_n)$, where $S_n=\lvert\mathcal{T}(G_n)\rvert$
  is the number of candidate partitions on $G_n$
  \citep[cf.][]{ChenSyrgkanisAustern2022}. Then \eqref{eq:En-rate} holds iff
  \begin{equation}\label{eq:recovery-rate}
    M_n\Delta_n-\log S_n\;\gg\;\tfrac12\log n.
  \end{equation}
  Under the adaptive log-$n$ grid, $\log S_n\asymp\log n$ (the grid has
  $\asymp(\log n)^{d}$ cells, so polynomially many candidate partitions) and
  $\Delta_n\asymp(\log n)^{-\kappa}$; hence \eqref{eq:recovery-rate} requires
  $M_n\gg(\log n)^{1+\kappa}$. \emph{Caveat (open):} the union bound assumes
  \emph{independent} split statistics over a \emph{fixed} candidate set $S_n$;
  the sequential/greedy optimal-tree search induces dependence and a data-grown
  $S_n$, so \eqref{eq:recovery-rate} is a sufficient condition for an idealized
  modal selector, not yet a theorem for the deployed selector
  (see \S\ref{sec:En-open}).
  \item[(A)] \emph{Held-out selection.} $\hat\tau$ is measurable with respect to
  auxiliary folds independent of the leaf-refit sample; then $\hat\tau$ is fixed
  conditional on that sample and no recovery event is needed
  ($E_n$ is the full space, $\PP(E_n^c)=0$, trivially \eqref{eq:En-rate}). The
  cost is that selection uses a separate fold, reintroducing splitting at the
  selection stage (not at inference).
\end{itemize}
\end{assumption}

\begin{theorem}[Working-model CLT]\label{thm:clt}
Suppose Assumptions~\ref{ass:positivity}--\ref{ass:En} hold and the tree nuisance
rates satisfy the CSA consistency condition
$\E_X\Ltwo{\Ghat(X)-\gz(X)}^2=\odiff_p(n^{-1/2})$
(equivalently each nuisance $\odiff_p(n^{-1/4})$; delivered by the tree oracle
rate when $\beta>d/2$). Then, on $E_n$,
\begin{equation}\label{eq:clt}
  \sqrt n\,\bigl(\htheta-\thstar\bigr)\;\xrightarrow{d}\;N\!\bigl(0,\,V\bigr),
  \qquad
  V=\E\bigl[\psiscore(O;\thstar,\etau,\mtau)^2\bigr],
\end{equation}
and a LOO-stability plug-in $\hat V\xrightarrow{p}V$; hence
$\htheta\pm z_{1-\alpha/2}\sqrt{\hat V/n}$ is an asymptotically valid CI for
$\thstar$. Because $\PP(E_n)\to1$, \eqref{eq:clt} holds unconditionally.
\end{theorem}

\begin{proof}
Neyman orthogonality of \eqref{eq:score} holds (Lemma~1, Remark on
orthogonality: both Gateaux partials at the truth vanish), and the second-order
smoothness $D_{gg}M=O(\Ltwo{g-\gz}^2)$ is the bilinear form \eqref{eq:lemma1}.
On $E_n$ the partition equals the fixed $\tau_0$, so
Lemma~\ref{lem:leafstab} supplies CSA (S3)--(S4)+MSC, hence stochastic
equicontinuity. The remaining CSA regularity holds: $A(\gz)=-1$ is invertible;
$A$ is $g$-free (trivially Lipschitz); $\mathrm{Var}(a(O;\gz))=\mathrm{Var}(A/\pi)<\infty$.
Then \citet[Theorem~1]{ChenSyrgkanisAustern2022} gives
$\sqrt n(\htheta-\theta(\etau,\mtau))\xrightarrow{d}N(0,A(\gz)^{-1}\Sigma_0 A(\gz)^{-1})$
with $A(\gz)^{-1}=-1$ and $\Sigma_0=\E[\psiscore^2]$, i.e.\ \eqref{eq:clt} with
target $\theta(\etau,\mtau)=\thstar$. Consistency of $\hat V$ follows from the
same stability (LOO variance estimator). Finally, we pass from the conditional
to the unconditional law. Decompose
$\sqrt n(\htheta-\thstar)=\sqrt n(\htheta-\thstar)\mathbf{1}_{E_n}
+\sqrt n(\htheta-\thstar)\mathbf{1}_{E_n^c}$. Under positivity the ATT is
bounded, so on $E_n^c$ the selected tree targets a different bounded working
model and $\htheta-\thstar=O_p(1)$; hence the second term has
$\E\bigl[\sqrt n\,\lvert\htheta-\thstar\rvert\mathbf{1}_{E_n^c}\bigr]
\lesssim \sqrt n\,\PP(E_n^c)=\odiff(1)$ by \eqref{eq:En-rate}, and it is
$\odiff_p(1)$. Equivalently, the total-variation distance between the conditional
law $\mathcal L(\sqrt n(\htheta-\thstar)\mid E_n)$ and the unconditional law is
$O(\PP(E_n^c))=\odiff(1)$, so the two limits coincide and \eqref{eq:clt} holds
unconditionally.
\end{proof}

\begin{remark}[Why $\odiff(n^{-1/2})$, not $\odiff(1)$, is necessary]\label{rem:rate-necessity}
The strengthening \eqref{eq:En-rate} over the naive $\PP(E_n)\to1$ is not
cosmetic. If $\PP(E_n^c)$ vanished only at rate $\odiff(1)$ but slower than
$n^{-1/2}$ (e.g.\ $\PP(E_n^c)\asymp n^{-1/4}$), the recovery-failure term would
contribute $\sqrt n\,\PP(E_n^c)\to\infty$ to the bias, breaking the CLT centering
even though structure recovery ``succeeds with probability tending to one.''
This is the post-selection subtlety: consistency of selection is insufficient;
the failure probability must beat the estimator's own $\sqrt n$ rate. It is also
why selection cannot be certified by simulation alone at moderate $n$---a $2$--$5\%$
empirical misrecovery rate looks benign but its \emph{rate in $n$} is what governs
validity.
\end{remark}

\section{The gap to \texorpdfstring{$\theta_0$}{theta0} and an impossibility}\label{sec:gap}

\begin{proposition}[Discretization gap]\label{prop:gap}
Under a H\"older-$\beta$ smoothness class for $(\ez,\mz)$ and the adaptive grid
$G_n$ ($b_n\asymp\log n$ bins per feature),
\[
  \thstar-\theta_0 = O\!\bigl((\log n)^{-2\beta}\bigr),
  \qquad\text{hence}\qquad
  \sqrt n\,\bigl(\thstar-\theta_0\bigr)\longrightarrow\infty
  \ \text{ for every fixed }\beta>0.
\]
\end{proposition}

\begin{proof}
By \eqref{eq:lemma1}, $\abs{\thstar-\theta_0}\le(\pi c)^{-1}
\Ltwo{\etau-\ez}\,\Ltwo{\mtau-\mz}$. On $G_n$ each grid-optimal
piecewise-constant nuisance has approximation error
$\Ltwo{\etau-\ez},\Ltwo{\mtau-\mz}=O(b_n^{-\beta})=O((\log n)^{-\beta})$ for a
H\"older-$\beta$ function on axis-aligned cells of width $1/b_n$ (discretization
note). The product is $O((\log n)^{-2\beta})$. Finally
$\sqrt n\,(\log n)^{-2\beta}\to\infty$ since any positive power of $n$ dominates
any power of $\log n$.
\end{proof}

\begin{corollary}[Impossibility for $\theta_0$]\label{cor:impossible}
No estimator built from a single \emph{fixed}, \emph{bounded-complexity}
($L=\odiff(\sqrt n)$) interpretable tree per nuisance on the adaptive grid admits
a $\sqrt n$-valid CI for $\theta_0$ under H\"older-$\beta$ smoothness. Driving the
gap to $\odiff(n^{-1/2})$ requires polynomial grid refinement, which forces
$L\not=\odiff(\sqrt n)$ (breaking Lemma~\ref{lem:leafstab}) or an off-tree
debiasing term (breaking ``the tree is the estimator'').
\end{corollary}

\begin{proof}
Immediate from Proposition~\ref{prop:gap}: the $\sqrt n$-scaled bias diverges, so
any CI of width $\Theta(n^{-1/2})$ centered at $\htheta$ (Theorem~\ref{thm:clt})
has $\theta_0$-coverage $\to0$. To make $\thstar-\theta_0=\odiff(n^{-1/2})$ one
needs $b_n^{-2\beta}=\odiff(n^{-1/2})$, i.e.\ $b_n\gtrsim n^{1/(4\beta)}$, giving
$L$ up to $b_n^{d}\gtrsim n^{d/(4\beta)}$; for this to also satisfy
$L=\odiff(\sqrt n)$ one needs $d/(4\beta)<1/2$, and even then the \emph{achieved}
complexity balancing estimation error is $s_n\asymp n^{d/(2\beta+d)}$, whose
approximation error $s_n^{-2\beta/d}\asymp n^{-2\beta/(2\beta+d)}$ is
$\odiff(n^{-1/2})$ only if $\beta>d/2$---but that regime is the cross-fit DML
result, not a fixed interpretable tree. The fixed-tree, log-grid object cannot
occupy it without a correction term absent from the displayed tree.
\end{proof}

\begin{corollary}[Structural PC special case: $\thstar=\theta_0$]\label{cor:structural}
If $(\ez,\mz)$ are \emph{exactly} piecewise-constant on some finite (unknown)
partition $\tau^\ast$ (a structural, not smoothness, assumption), then for $n$
large enough that $G_n$ resolves $\tau^\ast$ we have $\etau=\ez$, $\mtau=\mz$,
hence $\thstar=\theta_0$, and Theorem~\ref{thm:clt} delivers a $\sqrt n$-valid CI
for the \emph{true} ATT $\theta_0$.
\end{corollary}

\begin{proof}
Under exact representability the approximation errors in
Proposition~\ref{prop:gap} vanish once $\tau^\ast\subseteq G_n$; then
\eqref{eq:lemma1} gives $\thstar-\theta_0=0$ and Theorem~\ref{thm:clt} applies
verbatim with $\thstar=\theta_0$.
\end{proof}

\begin{remark}[The diagnostic $\hat\delta$ is not a correction]
Define $\hat\delta=\htheta-\htheta_{\mathrm{cf}}$, the difference between the
displayed-tree estimate and a cross-fit anchor. By \eqref{eq:lemma1} $\hat\delta$
is a plug-in estimate of the gap and is \emph{reported as an honesty flag}: it is
never subtracted from $\htheta$. Indeed, by Proposition~\ref{prop:gap} the gap is
asymptotically non-negligible relative to the CI width, so a $\hat\delta$-correction
capable of recovering $\theta_0$ would already constitute a solution of the
$\theta_0$ problem, contradicting Corollary~\ref{cor:impossible}.
\end{remark}

\section{Discharging the partition-fixing event \texorpdfstring{$E_n$}{En}: open problem}\label{sec:En-open}

Theorem~\ref{thm:clt} is stated \emph{conditional} on Assumption~\ref{ass:En},
i.e.\ on a selector delivering $\PP(E_n^c)=\odiff(n^{-1/2})$. Discharging that
assumption for the \emph{deployed} optimal-tree-over-grid selector is the one
remaining proof obligation. We record precisely what each route requires and
where it breaks, so the choice can be made once the empirical recovery rate is in
hand (the structure-recovery frequency $\widehat{\PP}(E_n)$ from the companion
simulation is a direct read on which route is viable).

\paragraph{Route B (exact recovery) --- what must be proved.}
Establish \eqref{eq:recovery-rate}, $M_n\Delta_n-\log S_n\gg\tfrac12\log n$, for
the actual selector. Two hazards, both flagged and neither yet closed:
\begin{enumerate}
  \item \emph{Dependence / data-grown candidate set.} The union bound
  $\PP(E_n^c)\le S_n\exp(-M_n\Delta_n)$ assumes independent split statistics over
  a fixed $S_n$. The greedy/optimal search is sequential (nested splits,
  discontinuous jumps) and $S_n=S_n(G_n)$ grows with the grid. A correct proof
  needs either (i) a concentration inequality for the modal winner under
  dependence (e.g.\ a bounded-differences / stability argument on the selection
  statistic), or (ii) a covering/peeling bound that controls the growing $S_n$
  without a naive union.
  \item \emph{Margin positivity under approximate PC.} $\Delta_n>0$ is generic
  only in the structural-PC regime (Corollary~\ref{cor:structural}); when the
  signal is merely H\"older, the risk landscape is flat ($O(h^2)$) near the
  optimum and $\Delta_n\downarrow0$ as the grid refines. Under the log-$n$ grid
  this is only $\Delta_n\asymp(\log n)^{-\kappa}$, which \eqref{eq:recovery-rate}
  can still tolerate---but the constant and the exponent $\kappa$ must be pinned,
  and the flat-landscape case is exactly where the empirical recovery rate is
  expected to degrade.
\end{enumerate}
\emph{Verdict:} viable iff the empirical recovery rate stays high as $n$ grows on
the target regimes; if it plateaus below one (as a flat margin predicts), Route B
cannot deliver $\odiff(n^{-1/2})$ and Route A is forced.

\paragraph{Route A (held-out selection) --- what must be proved.}
Trivially $\PP(E_n^c)=0$ (no recovery needed), so \eqref{eq:En-rate} is immediate
and the two hazards of Route B vanish. The obligations shift: (i) show the
held-out selection fold is asymptotically negligible for efficiency (the leaf
refit still uses the full inference sample, so the only cost is the selection
fold's information, $\odiff(1)$ of the total under a vanishing selection
fraction); (ii) state honestly that splitting re-enters at \emph{selection} (not
at inference), which weakens---but does not void---the ``no sample-splitting''
headline. This route is robust to selection instability by construction and does
\emph{not} require a margin.

\paragraph{Decision rule (keyed to the simulation).}
Let $r_n=\widehat{\PP}(E_n)$ be the measured recovery rate.
\begin{itemize}
  \item If $r_n\to1$ fast on the structural-PC (binary) regimes and stays high on
  approximate-PC as $n$ grows: pursue Route B; prove \eqref{eq:recovery-rate} via
  a dependence-robust concentration bound.
  \item If $r_n$ plateaus $<1$ (the flat-margin signature; the pilot's
  $\approx0.5$--$0.6$ on the complex DGP is a warning): adopt Route A as the
  primary discharge, cite Route B only under the structural-PC assumption where
  the margin is $O(1)$.
\end{itemize}
Either way the estimand and the CLT (Theorem~\ref{thm:clt}) are unchanged---only
the mechanism that supplies $E_n$ differs. The estimand fork ($\thstar$ vs
$\theta_0$) is orthogonal to the A-vs-B choice.

\paragraph{The empirical verdict (companion simulation).}
The structure-recovery frequency is stable across runs: $r_n=1.0$ on the
structural-PC (\texttt{simple}, \texttt{moderate}) regimes but $r_n\approx0.5$--$0.6$
on the \texttt{complex} regime, with no trend toward $1$ as $n$ grows over
$n\in\{500,\dots,10000\}$. This is the flat-margin signature: on the hard regime
$\PP(E_n^c)\approx0.4$ does not vanish, let alone at $\odiff(n^{-1/2})$. By the
decision rule, \textbf{Route~A is the primary discharge for the general case};
Route~B applies only under the structural-PC assumption
(Corollary~\ref{cor:structural}), where $r_n\to1$ is both predicted and observed.
We now make Route~A a theorem.

\begin{lemma}[Held-out selection discharges $E_n$]\label{lem:routeA}
Partition the sample into a selection fold $\mathcal S_n$ of size $\lfloor\rho_n
n\rfloor$ and an inference fold $\mathcal I_n$ of size $n-\lfloor\rho_n n\rfloor$,
with selection fraction $\rho_n\to0$ and $\rho_n n\to\infty$. Let $\hat\tau$ be
\emph{any} partition selector measurable with respect to $\mathcal S_n$ only, and
refit leaf values on $\mathcal I_n$. Then:
\begin{enumerate}
  \item[(i)] Conditional on $\mathcal S_n$, the partition $\hat\tau$ is fixed, so
  $E_n$ holds with $\PP(E_n^c\mid\mathcal S_n)=0$; \eqref{eq:En-rate} is satisfied
  trivially and Lemma~\ref{lem:leafstab} applies to the fixed $\hat\tau$ on
  $\mathcal I_n$.
  \item[(ii)] The estimator on $\mathcal I_n$ obeys
  $\sqrt{n-\lfloor\rho_n n\rfloor}\,(\htheta-\thstar)\xrightarrow{d}N(0,V)$, and
  since $\rho_n\to0$ the efficiency loss relative to the full sample is
  $\sqrt{n/(n-\lfloor\rho_n n\rfloor)}=1+\odiff(1)$; i.e.\ the asymptotic
  variance is unchanged and the CI attains nominal coverage for $\thstar$.
\end{enumerate}
\end{lemma}

\begin{proof}
(i) $\hat\tau$ is $\sigma(\mathcal S_n)$-measurable and $\mathcal S_n\indep
\mathcal I_n$, so conditional on $\mathcal S_n$ the partition is a deterministic
constant and the LOO perturbations of the $\mathcal I_n$ leaf means leave it
unchanged---exactly the fixed-partition hypothesis of Lemma~\ref{lem:leafstab},
whose $\nu$-stability rate $O(L/|\mathcal I_n|)=\odiff(|\mathcal I_n|^{-1/2})$
holds since $L=\odiff(\sqrt n)$ and $|\mathcal I_n|=n(1+\odiff(1))$.
(ii) Apply Theorem~\ref{thm:clt} on $\mathcal I_n$ (all its hypotheses hold there,
$\hat\tau$ now genuinely fixed), giving the CLT at rate
$\sqrt{|\mathcal I_n|}$. Because $|\mathcal I_n|=n-\lfloor\rho_n n\rfloor
=n(1-\rho_n)$ with $\rho_n\to0$, $\sqrt{n/|\mathcal I_n|}=(1-\rho_n)^{-1/2}\to1$,
so rescaling to the $\sqrt n$ rate leaves $V$ unchanged. Marginalizing over
$\mathcal S_n$ preserves the conclusion since it holds for every realization.
\end{proof}

\begin{remark}[Honest scope of ``no sample-splitting'']
Route~A reintroduces a split at the \emph{selection} stage---$\hat\tau$ is chosen
on $\mathcal S_n$, not on all $n$. What Route~A preserves is the property CSA
buys: \emph{no split at the inference stage}, i.e.\ leaf values and the CLT use
the full inference fold $\mathcal I_n$ with LOO stability standing in for
cross-fitting. This is strictly weaker than the ideal ``single tree, all $n$,
no split anywhere'' that Route~B targets under structural PC, but it is the price
of validity when the margin is flat and exact recovery fails---as the simulation
shows it does on non-tree signal. The displayed tree is still a single
interpretable tree; only its \emph{structure} was chosen on a subsample, while its
\emph{leaf values}---the numbers a reader acts on---use nearly all the data.
\end{remark}

\bibliographystyle{plainnat}
\begin{thebibliography}{1}
\bibitem[Chen et al.(2022)]{ChenSyrgkanisAustern2022}
Q.~Chen, V.~Syrgkanis, and M.~Austern.
\newblock Debiased machine learning without sample-splitting for stable
  estimators.
\newblock \emph{NeurIPS}, 2022. arXiv:2206.01825.
\end{thebibliography}

\end{document}
